\chapter{背景および関連研究}

\section{グリッチ表現}

グリッチ（glitch）は、一般に機器やシステムに生じる一時的な不具合や異常を指す語であり、特に電子機器やコンピュータなどの技術的な文脈で用いられてきた。デジタルメディアにおいては、データの記録、伝送、符号化・復号、描画などの過程において、本来意図されていない状態が生じることで、画像の欠損、色のずれ、画面の分断、ノイズなどとして現れる場合がある。Menkmanは、このようなグリッチを、技術的な処理の流れの中に生じる予期しない断絶として捉え、デジタルメディアが本来内包している不完全性の一つとして論じている\cite{menkman2011}。

一方、glitchという語は必ずしも視覚表現のみを指すものではない。現在では、コンピュータシステム上の一時的な不具合に加え、ゲームにおいて本来想定されていない挙動や現象を指して用いられることもある。このように、glitchという語が示す対象は文脈によって異なるが、いずれの場合にも、本来想定された正常な状態からの逸脱という性質を持つ。

また、グリッチと類似して用いられる語として、エラー（error）やノイズ（noise）がある。これらの語は相互に関連するものの、必ずしも同義ではない。本研究では、エラーをシステムや処理において生じる異常や失敗そのものを示す語として扱う。また、ノイズについては、本来の信号や情報に対して付加される不要な情報、あるいは視覚的には砂嵐や粒状の乱れなどとして現れるものを主に指す。これに対してグリッチは、こうしたエラーやノイズを含みながら、正常な表示や処理から逸脱した結果として現れる視覚的・時間的な変化を含む概念として扱う。

本来、グリッチは意図されない故障や異常によって生じるものである。しかし、その特徴的な視覚的性質は、次第に美的表現としても利用されるようになった。データの破損や伝送エラーなどによって偶発的に生じた画像を作品として利用するだけでなく、実際には故障していないシステム上で、色ずれ、画像の分断、欠損、ノイズなどの特徴を意図的に生成する表現も存在する。このように、グリッチそのもの、あるいはグリッチを想起させる視覚的特徴を利用した芸術表現は、一般にグリッチアートとして扱われている。

Moradiは、グリッチを用いた表現について、偶発的な故障やエラーによって生じたものを `Pure Glitch''、実際のグリッチに見られる視覚的特徴を意図的に再現したものを `Glitch-alike'' と分類している\cite{moradi2004}。Pure Glitchは、制作者によってあらかじめ設計されていない偶発的な現象であるのに対し、Glitch-alikeは、グリッチに特徴的な視覚表現を制作者が意図して生成したものである。したがって、グリッチを芸術表現として扱う際には、実際の故障によって生じたグリッチと、その外観や特徴を意図的に構成した表現とを区別して考える立場が存在する。

グリッチを芸術表現として扱う際には、意図せず発生したグリッチそのものと、グリッチに特徴的な外観を意図的に生成した表現とをどのように扱うかという問題が存在する。Moradiは、偶発的な故障やエラーによって生じたグリッチを `Pure Glitch''、グリッチの視覚的特徴を意図的に再現したものを `Glitch-alike'' と分類している\cite{moradi2004}。この分類が示すように、意図的に生成された表現をグリッチそのものとして扱うかについては、グリッチアートを論じる上で立場の違いが存在する。

そこで本研究では、「グリッチ」と「グリッチ表現」を区別して用いる。本研究における「グリッチ」は、情報の伝送、処理、記録、描画などの過程において意図せず発生する故障、異常、あるいは正常な状態からの逸脱を指す。一方、「グリッチ表現」は、グリッチによって生じる視覚的・動的特徴を利用した表現を指すものとする。

この「グリッチ表現」には、実際に生じたグリッチを素材として利用する場合だけでなく、色ずれ、画像の欠損、断片化、ノイズ、走査方向への変形など、グリッチを想起させる特徴を意図的に生成した表現も含める。これは、意図的に生成された現象がグリッチそのものであるか否かを判定するための定義ではなく、本研究が対象とする視覚表現の範囲を示すための呼称である。

本研究では、実際のシステム障害やデータ破損そのものを研究対象とするのではなく、これらによって生じる特徴的な外観や動きを視覚表現として利用することに着目する。したがって、以降では、作品やシステム上に提示されるこれらの視覚的・動的表現を総称して「グリッチ表現」と呼ぶ。



\section{デジタル表現としてのグリッチ}

前節で述べたように、グリッチは本来、情報処理や伝送などの過程において生じる異常な現象である。一方で、画像の断片化、色のずれ、走査方向への変形、ノイズなど、グリッチによって生じる視覚的特徴は、その特異な外観自体が美的表現としても利用されている。デジタルエラーを芸術的な手法として積極的に利用するグリッチアートについては、デジタルメディアに固有の不完全性や偶発性を表出する表現として、これまでにも作品制作および研究が行われている\cite{menkman2011,morgan2019}。

グリッチ表現の特徴の一つとして、デジタル媒体との強い結び付きが挙げられる。例えば、ブロック状の画像欠損は画像・映像圧縮を、RGB成分のずれはディスプレイ上の色表現を、走査方向に生じる帯状の変形や画面の断片化は映像信号やデータ処理の異常を想起させる。Morganは、デジタル画像に生じる圧縮アーティファクトが無秩序に見える一方で、実際には計算機の構造や圧縮方式によってその外観が規定されていることを指摘している\cite{morgan2019}。このように、グリッチ表現の外観は単なる抽象的な装飾ではなく、デジタルデータやそれを処理・表示する機構との関係を強く想起させる特徴を持つ。

また、グリッチ表現は、本来は表示内容を伝達するために存在するディスプレイやインタフェースそのものを、鑑賞者に意識させる表現となり得る。通常、映像や画像を閲覧する際、表示装置やデータ処理の存在はできる限り意識されず、その上に提示される内容が主な知覚対象となる。これに対し、画面全体に画像の欠損や色ずれ、走査線などが現れた場合、鑑賞者は表示されている内容だけでなく、「画面が正常に動作しているか」「データが破損しているのではないか」といった、媒体そのものの状態を意識する可能性がある。Menkmanは、こうしたノイズや破綻によって、通常は透明化されているメディアの構造や前提が表面化することをグリッチ表現の特徴として論じている\cite{menkman2011}。

この性質は、映像作品においても利用されている。例えば、映像そのものにグリッチを発生させるだけではなく、スマートフォンの通知やシステム画面を模した要素を画面上に配置するなど、鑑賞者が使用している表示装置やインタフェースと映像作品の境界を利用した演出が存在する。このような表現では、作品内部の出来事が実際の端末や画面上で発生しているかのように提示されることで、映像が表示される「画面」自体も表現の一部となる。グリッチ表現も同様に、表示内容のみではなく表示媒体そのものに異常が生じたかのような外観を作ることができるため、画面全体を利用した表現との親和性が高いと考えられる。

グリッチ表現は、静止画像と動画像の双方に用いることができる。静止画像では、ある時点における画像の欠損、色ずれ、断片化などを視覚的構成として提示することができる。一方、動画像では、これらの状態を時間的に変化させることが可能であり、画像の一部が瞬間的に移動する、ノイズが断続的に発生する、色成分のずれが変化するといった、異常が発生する過程そのものを表現として扱うことができる。静的なグリッチ表現と動的なグリッチ表現は視覚的特徴の多くを共有するが、後者では時間変化を表現要素として利用できる点に違いがある。

さらに、グリッチ表現はもともとデジタルデータの処理や変化と密接に関連するため、固定された映像として提示するだけでなく、外部から入力された情報に応じて表現を変化させることとも親和性がある。例えば、グリッチの発生位置、変形量、色ずれの大きさ、ノイズの発生頻度などを変数として扱うことで、時間経過やユーザ入力などに応じた動的な視覚表現として構成することが可能である。

本研究では、特にこのようなグリッチ表現の「画面そのものに異常が生じたように知覚される性質」と「時間的に変化させやすい性質」に着目する。これらの特徴を、スマートフォン上で日常的に表示され、画面全体を使用するライブ壁紙と組み合わせることで、映像作品として鑑賞する場合とは異なるグリッチ表現が可能になると考えた。


\section{動的・インタラクティブな視覚表現}

前節で述べたグリッチ表現は、静止画像として提示するだけでなく、時間経過に伴って状態を変化させる動的な表現として構成することができる。色ずれ、画像の断片化、帯状の変形、ノイズなどの視覚的特徴は、その発生位置や強度を時間的に変化させることで、静止状態とは異なる印象を与える。特に、グリッチは本来、正常な状態から突発的に逸脱する現象であるため、その発生や消失を時間変化として提示することは、グリッチに由来する不安定さや偶発性を表現する上で有効であると考えられる。

さらに、動的なグリッチ表現にユーザからの入力を取り入れることで、あらかじめ定められた映像を再生する場合とは異なる表現が可能となる。動画では、制作者が設定した時間軸に従って同一の変化が再生されるのに対し、インタラクティブな表現では、ユーザの行動や端末の状態に応じてその都度異なる変化を発生させることができる。これにより、表示されているグリッチを単なる映像上の演出としてではなく、現在使用している端末上で生じている現象のように知覚させることが可能になると考えた。

本研究では、ユーザ入力としてスマートフォン端末の傾きを利用する。タップやスワイプのような明示的な操作ではなく、端末を持つ、傾ける、向きを変えるといった通常の使用行動に伴って表現が変化する構成とした。このような入力方法では、ユーザが特定の位置を操作してエフェクトを発生させるのではなく、端末の状態変化に追従して画面全体の表現が連続的に変化する。そのため、操作対象としてのユーザインタフェースを強く意識させることなく、グリッチ表現にインタラクションを導入することができる。

また、傾きによる表示変化は、スマートフォンをはじめとする携帯端末における視差表現や動的な背景表現にも用いられている。本研究では、こうした端末の傾きに応じた表示変化をグリッチ表現へ応用する。例えば、端末の姿勢に応じて画像層の位置関係や色ずれ、変形量などを変化させることで、端末そのものを動かす行為と画面上の異常な変化とを結び付ける。

この構成には、実際の表示装置に生じる異常との視覚的な類似性もある。液晶ディスプレイでは、表示状態や故障の種類によっては、見る角度を変えることで色調や明度が変化したり、帯状の色変化が強く知覚されたりする場合がある。本研究では実際の故障状態を再現するものではないが、端末を傾けることで画面上の色や位置関係が変化する構成は、「表示装置そのものに異常が生じている」ような印象を形成する要素の一つとして利用できると考えた。

一方、本研究では、ユーザの操作とグリッチ表現との対応を過度に明示的にしないことも重視する。例えば、画面をタップした位置に必ずグリッチが発生する構成では、ユーザの入力と視覚的変化との因果関係が明確となり、グリッチよりも一般的な操作フィードバックとして認識される可能性がある。これに対して、端末姿勢の変化を入力とする場合、表現はユーザの行動に反応しながらも、グリッチが直接操作されているという印象を弱めることができる。本研究では、このような間接的なインタラクションによって、グリッチに由来する偶発性や不安定さを保ちながら、ユーザの行動に応じて変化する表現を目指す。

ただし、動的なグリッチ表現は、静的な背景と比較して視覚的な情報量が多く、ちらつきや急激な色・形状の変化によって、視認性の低下や視覚的負担を生じさせる可能性がある。特に、本研究で対象とするライブ壁紙は、一時的に鑑賞する映像作品とは異なり、日常的に繰り返し目にする可能性のある媒体である。したがって、表現の強さや変化量を単に大きくするのではなく、グリッチ表現としての特徴と、背景として利用する際の視認性や負担との関係を検討する必要がある。この点については、後述するライブ壁紙の特徴および評価実験において改めて扱う。




\section{ライブ壁紙}
\subsection{ライブ壁紙の概要}
スマートフォンにおける壁紙は、ホーム画面やロック画面の背景として表示される画像であり、端末の利用時に繰り返し目にする視覚要素の一つである。一般的な壁紙が静止画像を表示するものであるのに対し、ライブ壁紙は、時間経過やユーザ操作、端末状態などに応じて表示内容を動的に変化させる壁紙である。
ライブ壁紙は単に動画を背景として再生するものに限られない。例えば、時間帯に応じて景色を変化させるもの、画面へのタッチに応じてキャラクタや図形が反応するもの、端末を傾けることで背景に視差を生じさせるものなどが存在する。このようにライブ壁紙は、通常の壁紙が持つ「背景」という性質を維持しながら、アニメーションやインタラクションを導入できる表示形式である。
特にスマートフォンでは、ディスプレイ、タッチ入力、ジャイロセンサ、加速度センサなどが一つの端末に統合されているため、ライブ壁紙は端末の状態やユーザの身体的な操作を視覚表現へ直接反映することができる。そのため、ライブ壁紙は単なる装飾画像ではなく、端末の画面全体を利用した動的な視覚表現のための基盤として捉えることもできる。
\subsection{Androidにおけるライブ壁紙}
Androidでは、ライブ壁紙を実装するための仕組みとして\texttt{WallpaperService}が提供されている\cite{androidWallpaperService}。\texttt{WallpaperService}はAndroidのServiceを継承したコンポーネントであり、ライブ壁紙は独立した動画ファイルとしてではなく、Androidアプリケーションの一部として実装される。
ライブ壁紙の描画処理は\texttt{WallpaperService.Engine}を通して実行され、開発者はその内部で描画内容や更新処理を任意に実装できる。したがって、あらかじめ制作した映像を再生するだけでなく、時刻、ユーザ入力、センサ値、外部データなどを取得し、それらに応じてリアルタイムに描画内容を変化させることが可能である。
この点において、ライブ壁紙は一般的な動画壁紙とは性質が異なる。動画壁紙では、基本的に事前に生成された映像を時間軸に沿って再生するのに対し、プログラムとして実装されたライブ壁紙では、表示される内容そのものを実行時に生成・変更できる。例えば、ユーザが端末を傾けた方向に応じて背景画像の位置を変化させたり、タッチした位置に応じてオブジェクトを動かしたりすることができる。
また、Androidにおけるライブ壁紙の仕組みは比較的早期から提供されており、現在においてもOSの機能として維持・拡張されている。近年では、同一のライブ壁紙サービスから複数の異なる壁紙状態を扱う仕組みなども追加されており、動的かつユーザごとに異なる壁紙を扱うための基盤が整備されている\cite{android16Wallpaper}。
\subsection{ライブ壁紙の表現形態}
ライブ壁紙には多様な実装が存在するが、本研究では、表示内容を変化させる要因に着目し、大きく三つの形態に分類する。ただし、これらは相互に排他的なものではなく、複数の特徴を併せ持つものも存在する。
第一は、「映像再生型」である。これは、事前に制作された動画やアニメーションを壁紙として再生するものであり、時間経過に従って一定の映像が表示される。また、端末の傾きやホーム画面のスクロールに応じて映像の再生位置を変えるなど、一部の操作と組み合わせた例も存在する。
第二は、「状態連動型」である。これは、時間、天候、端末状態、ユーザの行動などに応じて、表示内容そのものを変化させるものである。例えば、時刻に応じて背景の景色を昼、夕方、夜へと変化させるものや、端末の使用状況を背景の状態として可視化するものなどが挙げられる。
第三は、「インタラクティブ型」である。これは、タッチ、スクロール、端末姿勢など、ユーザの操作や身体動作に応じてリアルタイムに表示内容を変化させるものである。画面上のキャラクタがタッチに反応する表現や、端末を傾けることで複数の画像レイヤに異なる移動量を与え、奥行きを感じさせる視差表現などがこれに該当する。
特にインタラクティブ型では、画面に表示された映像を一方向的に鑑賞するのではなく、実際に手に持っている端末の状態と視覚表現とが結び付く。例えば、端末を傾けることで画面内の空間の見える範囲が変化する表現では、ディスプレイが単なる平面的な表示面ではなく、その奥に存在する仮想的な空間を覗く窓のように知覚される場合がある。このように、ライブ壁紙はスマートフォンという物理的な端末と、画面内の視覚表現とを直接結び付けることのできる媒体である。
\subsection{スマートフォンOSにおける動的壁紙の現状}
動的な壁紙はAndroid以外のスマートフォンOSにも存在するが、その実装方法や利用可能なインタラクションはOSによって異なる。
Androidでは、前述した\texttt{WallpaperService}を利用することで、開発者が独自の描画処理を実装したライブ壁紙を提供できる。そのため、単純な動画再生に限らず、センサ入力、ユーザ操作、外部情報などを利用したインタラクティブな壁紙を構築することが可能である。
一方、iOSにおいても動的な壁紙表現は提供されており、Live Photosを利用したロック画面のアニメーションや、端末の動きに応じて奥行き表現を変化させる壁紙機能などが存在する\cite{appleLivePhotoWallpaper,appleSpatialScene}。ただし、Androidの\texttt{WallpaperService}のように、第三者のアプリケーションが任意の描画処理を継続的に実行するライブ壁紙とは実装方式が異なる。
このように、「動く壁紙」という概念自体は現在のスマートフォンにも広く存在する一方で、その自由度やインタラクションの方法はプラットフォームによって異なる。また、近年では写真や短い動画を利用した比較的単純な動的壁紙も容易に作成・利用できるようになっている。これに対して、本研究では、あらかじめ生成された映像を再生することではなく、端末の状態に応じてリアルタイムに描画内容そのものを変化させるインタラクティブなライブ壁紙に着目する。


\subsection{視覚表現媒体としてのライブ壁紙}
ライブ壁紙のもう一つの特徴は、独立したアプリケーションや映像作品ではなく、スマートフォンのユーザインタフェースの背景として存在することである。ユーザはライブ壁紙を鑑賞するために専用のアプリケーションを起動する必要がなく、ホーム画面を表示するという日常的な端末操作の中で、その表現を目にすることになる。
また、ライブ壁紙は画面の広い領域を利用しながら、アプリアイコン、時刻表示、ウィジェットなどのユーザインタフェースと同時に表示される。そのため、単独の映像作品とは異なり、「背景」と「インタフェース」の中間に位置する視覚表現として扱うことができる。
この性質は、本研究で扱うグリッチ表現と親和性が高い。グリッチ表現を映像作品の内部に配置した場合、鑑賞者はそれを作品中の演出として認識することができる。一方、スマートフォンの画面全体にグリッチ表現を提示し、その上に実際の時刻やアプリケーションアイコンなどが表示される場合、グリッチが作品内部だけでなく、端末の表示そのものに発生しているような印象を与えることができる。
さらに、端末の傾きなどに応じてグリッチ表現を変化させることで、ユーザが実際に手にしている端末の動作と、画面上の異常な変化とを関連付けることができる。本研究では、このような画面全体を利用できる点、実際のユーザインタフェースと共存できる点、およびユーザの操作に対してリアルタイムに反応できる点に着目し、ライブ壁紙を動的グリッチ表現の提示媒体として採用する。

% 図で静止壁紙／映像再生型／状態連動型／インタラクティブ型を4枚横に置いて、下に「固定」「時間」「外部状態」「ユーザ入力」と書く


\section{関連研究}

\subsection{グリッチ表現を用いた制作および研究}

グリッチを芸術表現として利用する試みについては、その美学的・文化的側面だけでなく、実際の制作過程に着目した研究も行われている。

Bedir Eriştiらは、新しいメディアアートとしてのグリッチについて、グリッチ表現の制作経験を持つデザイナによる作品制作と半構造化インタビューを実施し、グリッチを用いた制作過程および制作者による解釈について検討している\cite{bedireristi2019glitch}。この研究では、グリッチを単なる技術的なエラーとしてではなく、制作者が意図的に扱う視覚表現として捉え、その制作と美的な意味について考察している。

また、グリッチ表現には、色のずれ、画像の欠損、断片化、ノイズなど複数の視覚的要素が存在し、それらの組合せや強度によって大きく異なる外観を生成することができる。そのため、同一の画像を基にした場合でも、使用する処理やパラメータによって異なる表現を生成できる点に特徴がある。

このような、複数のパラメータから多様な視覚表現を生成できる性質は、グリッチに限らず生成的なデジタルアートにもみられる。Parikhは、311名の参加者がインタラクティブな生成アートを制作する際に選択した色や密度、線の特徴などを分析し、生成アート内部における複数の選好には一定の関係がみられることを報告している\cite{parikh2020preferences}。このことは、生成的な視覚表現において、単一の完成形を提示するだけでなく、ユーザ自身の選択を表現へ反映することに研究上の意義があることを示している。

さらに、インタラクティブアートにおけるユーザの関与について、作品内で操作可能なパラメータが多いほど、ユーザのエンゲージメントが高くなる可能性を示した研究も存在する\cite{interactiveArtEngagement}。このように、ユーザが表現へ介入し、その結果を自身の好みに合わせて変化させられることは、インタラクティブな視覚表現を設計する上で重要な要素の一つである。

本研究におけるグリッチ表現も、RGBずれ、画像の欠損、ノイズ、変形量など、複数の要素を独立したパラメータとして変更することができる。この性質を利用することで、一つの完成されたグリッチ表現を提示するだけでなく、利用者の視覚的な選好に応じて表現を変化させることが可能であると考える。


\subsection{ライブ壁紙を利用した情報提示}

ライブ壁紙を、単なる装飾的な背景ではなく、ユーザに情報を継続的に提示するインタフェースとして利用する研究も行われている。

MunsonらによるBeWell+では、ユーザの身体活動、睡眠、社会的交流などの状態を、魚や水中環境の変化としてスマートフォンのロック画面および壁紙に表示している\cite{bewellplus}。このシステムでは、ユーザが専用のアプリケーションを起動することなく、スマートフォンを確認する際に自身の状態を視覚的に把握できるよう、壁紙をambient displayとして利用している。

同様に、WhoIsZukiでは、ユーザの身体活動を物語の進行へ変換し、スマートフォンの壁紙上に提示することで、日常生活の中で身体活動を振り返るためのインタフェースを構築している\cite{whoiszuki}。このような研究では、壁紙がユーザの端末利用時に自然に視界へ入りやすいという特徴を利用し、情報を明示的に確認する操作を要求しない提示方法として活用されている。

Nwaguらは、無意識的なスマートフォン利用を抑制することを目的として、Androidライブ壁紙 ``Chai Wallpaper'' を提案している\cite{nwagu2023chai}。Chai Wallpaperでは、スマートフォンの利用頻度に応じて壁紙上の樹木の葉の色や状態を変化させ、ユーザ自身の利用状況を視覚的に提示している。その後の評価では、121名を対象として2週間の利用実験が行われ、無意識的なスマートフォン利用の減少が報告されている。また、壁紙を含むシステムのカスタマイズ機能についても、参加者から肯定的な評価が得られている。

これらの研究から、スマートフォンの壁紙は、ユーザが情報を見るための明示的な操作を行わなくても、端末を利用する過程で繰り返し情報を提示できる媒体として利用可能であることが分かる。特にライブ壁紙では、ユーザの行動や端末状態に応じて表示内容を継続的に変化させることで、静止壁紙とは異なる情報提示が可能となる。


\subsection{インタラクティブな視覚表現とユーザ選好}

インタラクティブアートでは、鑑賞者の入力が作品の一部として扱われ、入力に応じて視覚的・聴覚的な状態が変化する。ユーザによって操作可能な要素を設けることは、作品への関与や探索を促す要因となり得ることが報告されている\cite{interactiveArtEngagement}。

また、インタラクティブな対象に対する美的評価には個人差が存在することも示されている。したがって、動的な視覚表現を設計する場合、一つの表現をすべての利用者に対して最適なものとして提示するのではなく、複数の表現や調整可能なパラメータを用意し、ユーザごとの選好を反映できる構成を検討する余地がある。

本研究では、インタラクションそのものの新規性を目的とするのではなく、端末の傾きに応じた入力を、グリッチ表現の視覚的・動的特徴を強めるために利用する。また、グリッチの種類や強度を変更可能とすることで、インタラクティブなライブ壁紙に対するユーザごとの選好を反映できる構成を目指す。



\section{本研究の位置づけ}

これまで述べたように、グリッチおよびグリッチ表現については、美学的・文化的側面、制作手法、生成方法などさまざまな観点から研究が行われてきた。また、生成的・インタラクティブな視覚表現において、ユーザが複数のパラメータを操作することで多様な表現を生成する試みや、その過程におけるユーザの選好を扱った研究も存在する\cite{parikh2020preferences,interactiveArtEngagement}。

一方、ライブ壁紙に関する研究では、スマートフォンを利用する際に繰り返し目にするという壁紙の性質を利用し、ユーザの行動や状態を継続的に提示する試みが行われている。これらの研究では、ライブ壁紙を単なる装飾的な背景ではなく、ユーザと端末との間で情報を提示するインタフェースとして利用している。

このように、グリッチ表現、インタラクティブな視覚表現、ライブ壁紙についてはそれぞれ研究が行われている。しかし、グリッチ表現をインタラクティブなライブ壁紙として提示し、さらにその種類や強度をユーザの視覚的な選好に応じて調整することを一体的に扱った研究は限定的である。

本研究では、この点に着目し、グリッチ表現とインタラクティブなライブ壁紙を組み合わせるとともに、ユーザ自身が表現を調整可能なシステムを提案する。具体的には、RGBずれ、画像の欠損、ノイズ、帯状の変形などの複数のグリッチ表現をパラメータ化し、それぞれの強度や組合せを変更できるようにする。さらに、スマートフォンの傾きを入力として用いることで、ユーザの端末操作に応じてリアルタイムに変化する動的なグリッチ表現を実現する。

本研究では、すべてのユーザに共通して好まれる単一のグリッチ表現を決定することのみを目的とはしない。グリッチ表現に対する好みや、ライブ壁紙として利用したいと感じる表現には個人差が存在する可能性がある。そこで、複数の表現に対するユーザ評価を通して、グリッチ表現の種類や強度と選好との関係を調査するとともに、その個人差について検討する。

さらに、得られた結果をもとに、利用者ごとの選好に応じて表現を変更可能とすることの有効性を検討する。これにより、あらかじめ制作者が一つの完成形を定めるライブ壁紙ではなく、ユーザ自身が好みに応じて表現を調整できる、パーソナライズ可能なライブ壁紙の設計について考察する。

以上より、本研究は、グリッチ表現をインタラクティブなライブ壁紙上で提示することに加え、表現に対する選好の個人差を考慮した調整機能を導入し、その評価を通してライブ壁紙向け動的グリッチ表現の設計指針を得ることを目的とする。